\section{Governance: reproducibility and privacy}
\label{sec:governance}

A latent model estimated from human behavioural data must be reproducible enough
to audit and private enough to respect the people in it. The two requirements
pull in opposite directions---auditability wants to retain data, privacy wants to
delete it---and the design resolves them by separating what is anonymous from what
is personal.

\subsection{Reproducible by construction}

Each fit is pinned to a content hash of the exact cohort it consumed (a SHA-256 of
the canonicalised outcome table). A cited estimate---a difficulty function or a
rider's skill posterior from cohort hash $\mathtt{abc123\ldots}$---can be
re-derived by re-running the estimator on the cohort that hash references, without
the authors' cooperation. Combined with the deterministic marginal-maximum-
likelihood estimator (Section~\ref{sec:mml}), this makes every published number
independently checkable: same cohort hash and same seed yield the same posterior.
Storing fits append-only, keyed by cohort hash, additionally turns a rider's skill
estimates over successive cohorts into a trajectory---a signal for whether skill
estimates are stable or drift after a regime change.

\subsection{Privacy: the anonymous--personal split}

The two output classes have different privacy status, and the split is the crux.
\begin{itemize}
  \item The \textbf{difficulty atlas} $\hat\diff(\metric,\site)$ is a property of
    a \emph{place}, not a person. It carries no per-rider information and is
    therefore anonymous; it survives deletion requests, in the same way that
    aggregate statistics do.
  \item The \textbf{skill posteriors} $\hat\skill_\rider$ are per-person
    attributes and are treated as sensitive: on a rider's opt-out, every skill
    estimate for that rider is hard-deleted across all cohort versions. Deletion
    is a fan-out over the append-only history, not merely a removal of the latest
    value.
\end{itemize}
This separation is what lets the atlas remain a durable, publishable scientific
artifact while the personal layer remains fully erasable. An anonymised public
export of the atlas---generalised in space and lagged in time, released only for
cells with enough distinct contributors---follows the same governance as other
anonymised environmental exports and does not depend on retaining any individual's
data.

\subsection{Skill is an input, not an inference}

Finally, a boundary that is as much a governance choice as a modelling one:
skill-conditioned scoring (Section~\ref{sec:model}) takes the skill level as an
\emph{explicit, caller-provided} input. The score for ``expert'' conditions is the
expert curve; the system does not infer a person's skill from their identity in
order to score them. This keeps the objective, physics-anchored difficulty curve
separate from any automatic profiling of individuals---skill-conditioned physics,
not personalisation---and it means the difficulty atlas can be built and used
without ever attaching a skill estimate to a named person.
